\subsection{Entropy and energy}
It was proved in \cite{DGL20, DNL21} that \[\Ent(X,\theta)\cap \mathcal{E}(X,\theta) \subset \cE^{\frac{n}{n-1}}(X,\theta).\] As will be shown below, one can extend these results to the case of prescribed singularities.

We start with an integrability result of Moser--Trudinger type for general weights:

\begin{Theorem}
 	\label{thm: rel Moser--Trudinger inequality weight}
 	Let $\phi$ be a model potential with $m_\phi>0$. Let $\chi_1\colon \R^+\rightarrow \R^+ $ be a weight. Let \smash{$\chi_2(t) := \int_0^t \chi_1(s)^{\frac{1}{n}} {\rm d}s$}. Then there exist $c>0$, $C>0$ depending on $X$, $\theta$, $n$, $m_\phi$ and $\omega$ such that, for all $\varphi\in \mathcal{E}_{\chi_1}(X,\theta, \phi)$ with $\sup_X \varphi=-1$, we have
 	\begin{equation}
 	\int_X \exp \big( c E_{\chi_1}(\varphi, \phi)^{-\frac{1}{n}} \chi_2(\phi-\varphi)\big)\omega^n \leq C.
 	\label{eq: Moser} \end{equation}
 \end{Theorem}

 \begin{proof}
 We first note that if we replace $\chi_1$ by $\alpha \chi_1$ with $\alpha$ positive constant the left-hand side of~\eqref{eq: Moser} does not change, so we may assume $\chi_1(1) = 1$. 	
We claim that it suffices to prove the above inequality for $\varphi$ with relative minimal singularities, i.e., $[\varphi]=[\phi]$. Indeed, given~${\varphi\in \psh(X, \theta)}$, $\varphi_j:=\max(\varphi, \phi-j)$ has the same singularity type as $\phi$. This would mean that
 \[\int_X \exp \big( c E_{\chi_1}(\varphi_j, \phi)^{-\frac{1}{n}} \chi_2(\phi-\varphi_j)\big)\omega^n \leq C.\]
Moreover, it follows from Lemma \ref{energy increases} that $ E_{\chi_1}(\varphi_j, \phi) \nearrow E_{\chi_1}(\varphi, \phi) $ as $j\rightarrow +\infty$. Fatou's lemma will then give the desired estimate.

Thus assume that $\varphi$ has relative minimal singularities.
Let $\psi=-a\chi_2(\phi-\varphi)+\phi$ where $a>0$ is a small constant to be suitably chosen. Then define $u=P_\theta(\psi)$ and $v=-\gamma_2 (\phi-u) +\phi$, where~${\gamma_2 \colon \mathbb{R}^+ \rightarrow \mathbb{R}^+}$ denotes the inverse function of $a\chi_2$, which is concave and increasing. We observe that $u$ is a $\theta$-psh function with the same singularity type as $\phi$ and that
\[v=-\gamma_2 (\phi-u) +\phi \leq -\gamma_2(\phi-\psi) +\phi= -\gamma_2(a\chi_2(\phi-\varphi)) +\phi = \varphi\] with equality on the contact set $\mathcal{C}=\{u=\psi\}$.

A simple computation gives
 \begin{flalign*}
 \theta +{\rm dd}^c v & = \gamma_2'(\phi-u) \theta_u + (1-\gamma_2'(\phi-u)) \theta_{\phi} - \gamma_2''(\phi-u) d(\phi-u) \wedge d^c (\phi-u)\\
&\geq \gamma_2'(\phi-u) \theta_u+ (1-\gamma_2'(\phi-u)) \theta_{\phi},
 \end{flalign*}
where in the above we used that $\gamma_2$ is concave. We consider the set $G:=\{\gamma_2'(\phi-u)< 1\}$. We are going to show by contradiction that for a suitable choice of $a$ we have $G \neq X$. So, we assume $G= X$.
It then follow that $v$ is $\theta$-psh and, by construction, we infer that $v$ has the same singularity type as $\phi$.

Recall that $\sup_X \varphi=\sup_X(\varphi-\phi)$ as $\phi$ is a model potential \cite[Lemma 3.5]{DDL6}, hence $\phi-\varphi \geq 1$. This implies that $\chi_1$ is increasing and $\chi_1(1) = 1$ giving that
 \[E_{\chi_1}(\varphi, \phi) = \int_X \chi_1(\phi-\varphi) \theta_{\varphi}^n \geq \chi_1(1)\int_X \theta_{\varphi}^n= \int_X \theta_{\phi}^n=m_\phi.\]

By \cite[Theorem 2.7]{DDL6}, the non-pluripolar Monge--Amp\`ere measure $(\theta+{\rm dd}^c u)^n$ is supported on $\mathcal{C}$, hence
\begin{align*}%\label{eq: DGL contact weight}
	(\theta+{\rm dd}^c v) &\geq \gamma_2'(\phi-u) (\theta+{\rm dd}^c u)
		 = (a\chi_2'(\phi-\varphi))^{-1} (\theta+{\rm dd}^c u)\quad \text{on $\mathcal{C}$}.
\end{align*}
The last identity follows from the fact that, since $(a\chi_2)'(\gamma_2(t))\gamma_2'(t)=1$ and $\gamma_2(\phi-u)=\phi-\varphi$ in $\mathcal{C}$, we have $\gamma_2'(\phi-u)= (a\chi_2'(\phi-\varphi))^{-1}$ in $\mathcal{C}$.

It follows from \cite[Lemma 5.19]{DDL6} that the above inequality between positive currents implies an inequality between the non-pluripolar measures (observe that this is not trivial since $(a\chi_2'(\phi-\varphi))^{-1} (\theta+{\rm dd}^c u)$ is not closed).
Thus, we can infer that
\begin{equation*}
 {\bf 1}_{\mathcal{C}} \;(a\chi_2'(\phi-\varphi))^{-n} (\theta+{\rm dd}^c u)^n
	\leq {\bf 1}_{\mathcal{C}}(\theta+{\rm dd}^c v)^n \leq {\bf 1}_{\mathcal{C}}(\theta+{\rm dd}^c \varphi)^n,
\end{equation*}
where the last inequality follows from \cite[Lemma 4.5]{DDL5}. The above is equivalent to
\begin{align}
 (\theta+{\rm dd}^c u)^n &={\bf 1}_{ \mathcal{C}} (\theta+{\rm dd}^c u)^n \leq {\bf 1}_{ \mathcal{C}} (a\chi_2'(\phi-\varphi))^{n}(\theta+{\rm dd}^c \f)^n\nonumber\\
 &\leq a^n \chi_1(\phi-\varphi)(\theta+{\rm dd}^c \f)^n, \label{eq: stimabase}
\end{align}
where in the last inequality we used that $(\chi_2')^n=\chi_1$.

We now choose $a$ so that
\[
\beta a^n \int_X \chi_1(\phi-\varphi) (\theta+{\rm dd}^c \varphi)^n=\beta a^n E_{\chi_1}(\varphi, \phi) = m_\phi
\]
 with $\beta > 1$.
Observe that $a\in(0,1)$ since we observed that $E_{\chi_1}(\varphi, \phi)\geq m_\phi$.
Integrating both sides of \eqref{eq: stimabase} over $X$, we obtain
\[
\int_X (\theta+{\rm dd}^c u)^n \leq \frac{m_\phi}{\beta}.
\]
Since $\int_X \theta_u^n=m_\phi$, we arrive at a contradiction.

So we can infer that $G\neq X$, or equivalently that $G^c\neq \varnothing$. This means that there exists~${x_0\in X}$ such that $(\phi-u)(x_0)\leq \tau_2(1)$ where $\tau_2$ is the inverse function of $\gamma_2'$ which we note to be decreasing (since so in $\gamma_2'$). In particular,
\[\sup_X u =\sup_X (u-\phi) \geq -\tau_2(1).\]

Applying the uniform version of Skoda's integrability theorem \cite{Zer01} to $\PSH(X,C\omega)$ for ${C>0}$ such that $\theta\leq C\omega$, we know that there exist uniform constants $c_0, C_0>0$ such that, for all~${h \in \PSH(X,\theta)}$ with $\sup_X h=0$ we have $\int_X {\rm e}^{-c_0 h} \omega^n \leq C_0$. For $h:=u+\tau_2(1)$, we have $\sup_X h \geq 0$, hence
\begin{align*}%\label{key estimate}
 \int_X {\rm e}^{c_0 (a \chi_2(\phi-\varphi) -\tau_2(1))} \omega^n &= \int_X {\rm e}^{c_0 (-\psi+\phi -\tau_2(1))} \omega^n \leq \int_X {\rm e}^{-c_0 (u +\tau_2(1))} \omega^n
\\&= \int_X {\rm e}^{-c_0 (h-\sup_X h) - c_0\sup_X h} \leq C_0,
\end{align*}
where in the first inequality we used $\phi\leq 0$ and $u\leq \psi$.

Let us now give a more explicit expression of $\tau_2(1)$.
By definition, $\tau_2$ is such that
\[s=\tau_2 \left( \frac{1}{a\chi_2'(\gamma_2(s))} \right).\] Now we want to understand $\tau_2(t)$, hence we want to find $s$ such that \smash{$\tfrac{1}{a\chi_2'(\gamma_2(s))}=t $}. This means~${ (at)^{-1}=\chi_2'(\gamma_2(s)) = (\chi_1(\gamma_2(s)))^{1/n} }$ (since $(\chi_2')^n=\chi_1$). Therefore, $\gamma_1((at)^{-n}) =\gamma_2(s)$, where $\gamma_1$ is the inverse function of $\chi_1$. It the follows that $ \tau_2(t)= a\chi_2(\gamma_1((at)^{-n}))$.
In particular, $\tau_2(1)= a\chi_2(\gamma_1(a^{-n}))$.

Also, letting $s:=\gamma_1 (a^{-n})$, we have $a^{-1} = (\chi_1(s))^{\frac{1}{n}}$ so that $a=a(s) = \frac{1}{\chi_2'(s)}$ and
\begin{gather} \tau_2(1) = \frac{\chi_2(s)}{\chi_2'(s)}. \label{Bo} \end{gather}
Note that since $a\in(0,1)$, then $s > 1$ ($\gamma_1$ is increasing and $\gamma(1)=1$).

Next, setting \[K:= \{ x \in X \colon a \chi_2(\phi - \varphi) \leq 2(\phi - \varphi) \},\]
we claim that
\[2^{-1} a \chi_2(\phi - \varphi) \leq a \chi_2(\phi - \varphi) - \tau_2(1) \quad \text{on $X \setminus K$},\]
that is, $a \chi_2(\phi - \varphi) \geq 2 \tau_2(1)$ on $X \setminus K$.

We observe that, since $\chi_2$ is convex and $\chi_2(0) = 0$, the set $E:=\{ t \in \mathbb{R}^+ \colon a\chi_2(t) \leq 2 t \}$ is a~closed convex set containing $0$, i.e., it is an interval of the form $[0, \lambda]$ with $0 \leq \lambda \leq + \infty$. Hence, $x\in X\setminus K$ if and only if $(\phi-\varphi)>\lambda$.
We then need to prove that if $ a\chi_2( \phi - \varphi) > a\chi_2( \lambda)$,
then
$ a \chi_2( \phi - \varphi) \geq {2 \tau_2(1)}$. It is then sufficient to prove that $a\chi_2( \lambda) \geq 2 \tau_2(1)$. The above is equivalent to $ \lambda \geq \gamma_2(2 \tau_2(1)) \geq 0$. By definition of $\lambda$, this means that $\gamma_2(2 \tau_2(1))\in E$, i.e.,
\[\tau_2(1) \leq \gamma_2(2 \tau_2(1)),\qquad \text{or} \qquad a \chi_2(\tau_2(1)) \leq 2 \tau_2(1).\]
By (\ref{Bo}) and $a=1/\chi_2'(s)$, the above inequality gives
\[\chi_2\left(\frac{\chi_2( s)}{\chi_2' (s)} \right) \leq 2 \chi_2(s).\]
Since $\chi_2$ is convex and $\chi_2(0) = 0$, we have
$\chi_2(s) \leq s \chi_2'(s)$.
Then we obtain
\[ \chi_2\left(\frac{\chi_2( s)}{\chi_2' (s)} \right) \leq \chi_2(s) \leq 2 \chi_2(s),\]
and the claim is proved.

It follows that
\begin{align*}%\label{final}
\int_X {\rm e}^{\frac{c_0}{2} a \chi_2(\phi-\varphi)} \omega^n &\leq \int_K {\rm e}^{c_0 (\phi-\varphi)} \omega^n + \int_{X\setminus K}{\rm e}^{c_0(a\chi_2(\phi-\varphi) -\tau_2(1))} \omega^n\\
 &\leq \int_K {\rm e}^{-c_0 \varphi} \omega^n + \int_{X\setminus K}{\rm e}^{c_0(a\chi_2(\phi-\varphi) -\tau_2(1))} \omega^n \\
 &\leq\int_K {\rm e}^{-c_0 (\varphi+1)+c_0} \omega^n + C_0\leq C_0({\rm e}^{c_0}+1),
\end{align*}
wherein the last inequality we applied the uniform Skoda theorem to the function $\varphi + 1$. Recalling that $\beta a^n E_{\chi_1}(\varphi, \phi) = m_\phi$, the result then follows with
$c(\beta) = \frac{c_0}{2} \beta^{(\frac{-1}{n})}m_\phi^{1/n}$, and $C= C_0({\rm e}^{c_0}+1)$.
As $\beta>1$ can be chosen arbitrarily, we observe that since $\phi$ and $\varphi$ are assumed to be equisingular, the function $\exp \big( c(\beta) E_{\chi_1}(\varphi, \phi)^{-\frac{1}{n}} \chi_2(\phi-\varphi)\big)$ is uniformly bounded, so by Lebesgue dominated convergence theorem, we can choose \smash{$c = \frac{c_0}{2} m_\phi^{1/n}$}.
Observe that the constants $c$ and $C$ are independent of $\chi_1$ (and $\chi_2$).
\end{proof}
